% Generated draft: requires mathematical and source review.
\documentclass[11pt]{article}
\usepackage[margin=25mm]{geometry}
\usepackage{fontspec}
\setmainfont{Latin Modern Roman}
\usepackage{amsmath,amssymb,mathtools,unicode-math}
\setmathfont{Latin Modern Math}
\usepackage{enumitem,needspace,xurl,hyperref}
\setlength{\parindent}{0pt}
\setlength{\parskip}{5pt}
\setlength{\emergencystretch}{3em}
\newcommand{\N}{\mathbb{N}}
\newcommand{\Z}{\mathbb{Z}}
\newcommand{\Q}{\mathbb{Q}}
\newcommand{\R}{\mathbb{R}}
\newcommand{\C}{\mathbb{C}}
\newcommand{\F}{\mathbb{F}}
\newcommand{\A}{\mathbb{A}}
\newcommand{\PP}{\mathbb{P}}
\newcommand{\E}{\mathbb{E}}
\newcommand{\eps}{\varepsilon}
\newcommand{\dd}{\,\mathrm{d}}
\begin{document}
\section*{UnsolvedMath ID: 800777}
\textbf{Source problem code:} catalogue entry
\subsection{Definitions and mathematical statement}
Let \((X,\Delta)\) be a pair, where \(\Delta\) denotes its boundary divisor and \(K_X\) denotes a canonical divisor. The condition that the pair is Kawamata log terminal (klt) means that \(K_X+\Delta\) is \(\R\)-Cartier and every divisorial discrepancy is strictly greater than \(-1\). A divisor \(D\) is nef if \(D\cdot C\geq 0\) for every complete integral curve \(C\subset X\). The supplied excerpt does not explicitly specify the ground field, projectivity hypothesis, or whether boundary coefficients are rational or real; these conventions require confirmation from the cited source.

The maximal rationally connected (MRC) fibration depends on \(X\), whereas the slope rationally connected quotient (sRC quotient), attributed in the supplied record to [Cam16], incorporates the boundary \(\Delta\); the record also refers to [Cam04, Cam16]. The excerpt does not supply the definitions or conventions for orbifold morphisms, orbifold fibers, or slope rational connectedness. These terms below retain their source-specific meanings, which must be recovered from those references before the statement is self-contained.

The supplied Conjecture 12, attributed there to [CCM19, Conjecture 1.5], asserts the following. Let \((X,\Delta)\) be a klt pair such that \(X\) is smooth and \(-(K_X+\Delta)\) is nef. Then there exists an orbifold morphism, described as a fibration,
\[
\rho:(X,\Delta)\longrightarrow(R,\Delta_R)
\]
with all of the following properties:
\begin{enumerate}
\item The pair \((R,\Delta_R)\) is klt, the variety \(R\) is smooth, and
\[
c_1(K_R+\Delta_R)=0.
\]
The excerpt does not specify the cohomology group or equivalence convention intended by this equality.
\item For a general point \(r\in R\), the general orbifold fiber \((X_r,\Delta_r)\), with \(X_r=\rho^{-1}(r)\), is slope rationally connected.
\item The fibration is locally trivial with respect to pairs: for every sufficiently small open set \(U\subset R\), there is an isomorphism over \(U\)
\[
\bigl(\rho^{-1}(U),\Delta|_{\rho^{-1}(U)}\bigr)
\cong
\bigl(U,\Delta_R|_U\bigr)\times\bigl(X_r,\Delta_{X_r}\bigr),
\]
where \(X_r\) is a general fiber. Under the usual product-pair notation, the boundary on the right is
\[
p_U^*(\Delta_R|_U)+p_F^*\Delta_{X_r},
\]
where \(p_U\) and \(p_F\) are the two projections.
\end{enumerate}
The source uses both \(\Delta_r\) and \(\Delta_{X_r}\) for the fiber boundary without explaining their relationship. It also says that the displayed isomorphism is over \(U\subset Y\), although the stated base is \(R\) and \(Y\) is undefined. The display above reads the isomorphism as being over \(U\subset R\); this is an explicit editorial interpretation requiring review. The topology in which \(U\) is open is not specified.

The supplied narrative reports that the conjecture was solved for log smooth surface pairs in [CCM19, Theorem 1.6] and remained open in general. These are supplied source-status claims, not independently verified mathematical or literature assessments.
\subsection{Short English statement}
The conjecture predicts that a smooth klt pair with nef anticanonical log divisor admits an orbifold fibration onto a smooth klt base with vanishing log first Chern class, with slope rationally connected general orbifold fibers and local product decompositions respecting the boundaries. The excerpt leaves foundational conventions and some notation unresolved. Its reported surface result and general open status are supplied claims requiring review.
\subsection{Sources}
\begin{itemize}
\item Local UnsolvedMath registry, permanent ID \(800777\), exact source problem code \texttt{catalogue entry}, title ``Positively Curved Projective Varieties --- Conjecture 12.'' The record attributes the problem to Matsumura, with proposed year 2022, and labels it open. This is a supplied registry label, not an independently verified status.
\item Matsumura, \emph{Open Problems on Structure of Positively Curved Projective Varieties} (2022), Conjecture 12. Supplied source URL: \url{https://arxiv.org/abs/2211.09153}. The record gives access date 2026-10-16 and extraction method \texttt{arxiv\_latex}; these are supplied provenance statements, not a claim of access or verification in this draft.
\item [CCM19, Conjecture 1.5], identified by the supplied excerpt as the attribution for the conjecture. The excerpt also cites [CCM19, Theorem 1.6] for the claimed log smooth surface case. Full bibliographic details and a URL for [CCM19] were not supplied.
\item [Cam16], credited in the supplied excerpt with introducing the slope rationally connected quotient, and [Cam04, Cam16], cited there for its relationship to MRC fibrations and boundary data. Full bibliographic details and URLs were not supplied.
\end{itemize}
\end{document}
